\section{Welfare and Policy Design}\label{sec:welfare_policy}

The price decomposition has a direct welfare interpretation. A set-aside raises
government outlays because it changes the order statistic that determines
payment, but the social cost is not the price effect itself. Part of the extra
payment is a transfer to the protected winner; part reflects a real production
inefficiency from assigning the contract to a higher-cost firm; and part is the
tax distortion created by financing the additional public outlay. The policy
question is therefore sharper than whether SME support is desirable. It is
whether the incremental SME surplus generated by excluding rival bidders is
valuable enough to justify the allocative and fiscal cost of removing them from
the auction.

\subsection{Welfare accounting}

Let $p^V$ denote the expected procurement price under policy $V$, and let
$c_{(1)}^V$ denote the winner's production cost. Relative to the open benchmark
$S_1$, the per-auction welfare loss from the SME-only set-aside $V_0$ is
\begin{equation}
    \mathrm{Loss}(V_0)
    =
    \underbrace{c_{(1)}^{V_0} - c_{(1)}^{S_1}}_{
        \mathrm{DWL}_{\mathrm{alloc}}}
    +
    \underbrace{\lambda\left(p^{V_0} - p^{S_1}\right)}_{
        \mathrm{MCPF\ distortion}} .
    \label{eq:welfare_v8}
\end{equation}
The first term is allocative deadweight loss: the additional production cost
created when the restricted auction selects an SME whose cost exceeds the
lowest-cost firm in the full pool. The second term prices the marginal cost of
public funds; the baseline uses $\lambda=\lambdaMain$, consistent with the
standard public-finance range used to evaluate tax-financed spending
\citep{ballard1985,hendren2020,finkelstein2020}. The remaining component of
the government's extra payment is not destruction in this accounting. It is an
implicit transfer to SME bidders, and it becomes a social benefit only to the
extent that the planner places extra welfare weight on SME surplus.

\subsection{The welfare cost of exclusion}

Table~\ref{tab:welfare_policy_v8} summarizes the welfare accounting and the
non-exclusionary policy benchmark. In non-pharmaceuticals, the full set-aside
raises government payments by \welfDeltaGovNp{} of the reference price. Of this,
\welfDwlAllocNp{} is real allocative waste and \welfMcpfDistLthirtyNp{} is the
MCPF distortion at $\lambda=\lambdaMain$, yielding a welfare loss equal to
\welfLossPctLthirtyNp~percent of the open-regime price. In pharmaceuticals, the
same calculation is larger: the payment increase is \welfDeltaGovPh, the
allocative loss is \welfDwlAllocPh, the MCPF distortion is
\welfMcpfDistLthirtyPh, and the total welfare loss is
\welfLossPctLthirtyPh~percent of the open-regime price. The lower endpoints of
the bootstrap intervals remain economically large
(\welfLossPctLoNp~percent in non-pharmaceuticals and \welfLossPctLoPh~percent
in pharmaceuticals), so the welfare conclusion is not a small-sample artifact.

\begin{table}[!htbp]
\centering
\begin{threeparttable}
\caption{Welfare cost of exclusion and the price-preference benchmark}
\label{tab:welfare_policy_v8}
\small
\setlength{\tabcolsep}{5pt}
\begin{tabular}{lcccccc}
\toprule
& \multicolumn{4}{c}{Full SME set-aside $V_0$}
& \multicolumn{2}{c}{10\% preference $V_3$} \\
\cmidrule(lr){2-5}\cmidrule(lr){6-7}
Class
& $\Delta_{\mathrm{gov}}$
& DWL$_{\mathrm{alloc}}$
& MCPF
& Loss / $p_{S_1}$
& $\Delta p$
& SME win gain \\
\midrule
Non-pharma
& \welfDeltaGovNp
& \welfDwlAllocNp
& \welfMcpfDistLthirtyNp
& \welfLossPctLthirtyNp\%
& \optPrefShiftNp
& \prefSmeWinGainTenNp~pp \\
Pharma
& \welfDeltaGovPh
& \welfDwlAllocPh
& \welfMcpfDistLthirtyPh
& \welfLossPctLthirtyPh\%
& \vThreeShiftPh
& \prefSmeWinGainTenPh~pp \\
\bottomrule
\end{tabular}
\begin{tablenotes}\footnotesize
\item $V_0$ is the full SME-only set-aside under the observed post-policy SME
pool. $\Delta_{\mathrm{gov}}$, DWL$_{\mathrm{alloc}}$, MCPF, and $\Delta p$ are
reported in reference-price units. MCPF uses $\lambda=\lambdaMain$. $V_3$ is an
open auction with a 10 percent SME price preference for winner selection; all
firms remain eligible, and the government pays the actual winning bid. SME win
gain is relative to the open benchmark $S_1$.
\end{tablenotes}
\end{threeparttable}
\end{table}

The magnitudes are large enough to matter outside the simulated auction. At the
observed \adherenceGsixtyfivePost~percent Group-65 adherence rate, the
per-auction losses imply an annual welfare cost of roughly
R\$\,\welfAnnualBRLLo~million in S\~ao Paulo Group-65 procurement alone; across
the realistic adherence range of \welfAdherenceLoPct--\welfAdherenceHiPct
percent, the range is R\$\,\welfRangeLoBRL--\welfRangeHiBRL~million. These
figures are not the central identification result, but they discipline the
policy relevance of the mechanism: exclusion is not a symbolic preference. It
is a costly way to transfer surplus.

\subsection{The policy frontier}

The clean benchmark is a price preference rather than no SME support. Under
$V_3$, SMEs receive a 10 percent scoring advantage for winner selection, but
non-SMEs continue to participate and continue to discipline the price-forming
order statistic. The simulated price effect is essentially zero:
\optPrefShiftNp{} of the reference price in non-pharmaceuticals and
\vThreeShiftPh{} in pharmaceuticals. The instrument does not replicate the full
redistribution of $V_0$; the SME win-rate rises by
\prefSmeWinGainTenNp/\prefSmeWinGainTenPh{} percentage points rather than being
forced to 100 percent. That difference is exactly the trade-off the set-aside
hides. Price preferences buy some reallocation toward SMEs while preserving the
competitors that set the auction's competitive outside option; set-asides buy
more reallocation by removing that outside option.

This frontier is particularly informative because the low-cost region is not a
knife-edge at 10 percent. In the preference grid, welfare costs remain inside
Monte Carlo noise up to \prefSafeBandHiNp{} percent in non-pharmaceuticals and
\prefSafeBandHiPh{} percent in pharmaceuticals, and even a 30 percent
preference costs only \welfPrefThirtyNp/\welfPrefThirtyPh{} percent of
$p_{S_1}$. The set-aside is therefore not the natural default instrument in
thick standardized procurement. It is an extreme point on the policy frontier:
maximum formal eligibility protection, purchased by deleting rival bidders from
the auction.

\subsection{How much must the planner value SME surplus?}

The distributional case for exclusion can be stated as a welfare-weight
condition. Following the social marginal welfare weight logic in
\citet{saezstantcheva2016}, let $w^{\mathrm{SME}}$ be the planner's weight on
one real of SME producer surplus relative to one real of general surplus. The
planner is indifferent between the set-aside and the 10 percent preference when
\begin{equation}
    w^{\mathrm{SME}}_\star
    =
    \frac{\mathrm{Loss}(V_0)-\mathrm{Loss}(V_3)}
         {\mathrm{SMESurplus}(V_0)-\mathrm{SMESurplus}(V_3)} .
    \label{eq:welfare_weight_v8}
\end{equation}
Because $V_3$ has near-zero welfare cost, the numerator is essentially the
deadweight cost of exclusion. The denominator is the additional SME surplus
created by moving from a price preference to a full set-aside. Under the main
specification, the implied indifference weights are
\welfWeightMainNp{} in non-pharmaceuticals and \welfWeightMainPh{} in
pharmaceuticals. In words, the planner must value an extra real of SME surplus
at more than twice a real of general surplus to prefer excluding rival bidders
over preserving them with a price preference.

The strict-invariance benchmark clarifies the scope of that statement. In
non-pharmaceuticals, the preference rule continues to dominate the set-aside;
the implied weight remains above the utilitarian benchmark. In pharmaceuticals,
the ranking is sensitive to whether the post-policy SME pool is interpreted as
a selected equilibrium participant pool or forced to have the same primitive
cost distribution as the pre-policy pool; under strict invariance the implied
weight falls to \welfWeightSiPh. That sensitivity is not a nuisance detail. It
is the policy boundary of the paper: the argument against bidder exclusion is
strongest where the protected pool is thick enough that redistribution can be
achieved without changing the price-forming bidders.

\subsection{When is exclusion defensible?}

The results do not imply that set-asides are never defensible. They imply that
the defense must be explicit. Exclusion becomes more plausible when the
protected pool is thin, when the planner's objective is to create protected
market access rather than merely tilt marginal awards, when dynamic SME
capacity gains are expected to be large, or when the planner places a high
social weight on SME surplus. Those conditions are not identified by the legal
threshold alone. They are market-structure conditions.

Non-pharmaceutical Group-65 procurement is the hard case for the set-aside: the
SME pool is thick, goods are standardized, and a price preference preserves
competition while delivering SME wins at almost no welfare cost. Pharmaceuticals
are the boundary case: the SME pool is thinner, composition changes more under
the policy, and the welfare ranking depends on how the post-policy SME
distribution is interpreted. The broader design principle is therefore
conditional rather than universal. Use exclusion only when the planner is
prepared to pay for the missing competitors; otherwise, support SMEs with an
instrument that keeps those competitors in the auction.
